\documentclass[14pt,a4paper]{scrreprt}

\usepackage{cmap}
\usepackage[T1]{fontenc} 
\usepackage[utf8]{inputenc}
\usepackage[english,russian]{babel}

\usepackage{caption}
\usepackage{subcaption}

\usepackage{float}

\usepackage{enumitem}

\usepackage{graphicx}
\usepackage{multirow}


\usepackage{pgfplots}
\pgfplotsset{compat=newest}
\usepgfplotslibrary{units}

\usepackage{longtable}

\usepackage{caption}
\captionsetup{labelsep=endash}
\captionsetup[figure]{name={Рисунок}}
\captionsetup[subtable]{labelformat=simple}
\captionsetup[subfigure]{labelformat=simple}
\renewcommand{\thesubtable}{\text{Таблица }\arabic{chapter}\text{.}\arabic{table}\text{.}\arabic{subtable}\text{ --}}
\renewcommand{\thesubfigure}{\text{Рисунок }\arabic{chapter}\text{.}\arabic{figure}\text{.}\arabic{subfigure}\text{ --}}


\usepackage{textcomp}

\usepackage{amsmath}
\usepackage{amsfonts}
\usepackage{array}

\usepackage{geometry}
\geometry{left=30mm}
\geometry{right=15mm}
\geometry{top=20mm}
\geometry{bottom=20mm}
\geometry{foot=1.7cm}

\usepackage{titlesec}
\titleformat{\section}
{\normalsize\bfseries}
{\thesection}
{1em}{}
\titlespacing*{\chapter}{0pt}{-30pt}{8pt}
\titlespacing*{\section}{\parindent}{*4}{*4}
\titlespacing*{\subsection}{\parindent}{*4}{*4}

% Маркировка для списков
\def\labelitemi{$\circ$}
\def\labelitemii{$*$}


\usepackage{setspace}
\onehalfspacing % Полуторный интервал

\frenchspacing
\usepackage{indentfirst} % Красная строка

\usepackage{titlesec}
\usepackage{xcolor}
% Названия глав
\titleformat{\section}{\normalsize\textmd}{\thesection}{1em}{}

\definecolor{gray35}{gray}{0.35}

\newcommand{\hsp}{\hspace{20pt}} % длина линии в 20pt

\titleformat{\chapter}[hang]{\Huge}{\textcolor{gray35}{\thechapter.}\hsp}{0pt}{\Huge\textmd}

\titleformat{\section}{\Large}{\textcolor{gray35}\thesection}{20pt}{\Large\textmd}
\titleformat{\subsection}{\Large}{\thesubsection}{20pt}{\Large\textmd}
\titleformat{\subsubsection}{\normalfont\textmd}{}{0pt}{}

% Настройки введения

\addtocontents{toc}{\setcounter{tocdepth}{2}}
\addtocontents{toc}{\setcounter{secnumdepth}{1}}

\usepackage{tocloft,lipsum,pgffor}

\addtocontents{toc}{~\hfill\textnormal{Страница}\par}

\renewcommand{\cftpartfont}{\normalfont\textmd}

\addto\captionsrussian{\renewcommand{\contentsname}{Содержание}}
\renewcommand{\cfttoctitlefont}{\Huge\textmd}

\renewcommand{\cftchapfont}{\normalfont\normalsize}
\renewcommand{\cftsecfont}{\normalfont\normalsize}
\renewcommand{\cftsubsecfont}{\normalfont\normalsize}
\renewcommand{\cftsubsubsecfont}{\normalfont\normalsize}

\renewcommand{\cftchapleader}{\cftdotfill{\cftdotsep}}

\usepackage{listings}
\usepackage{xcolor}

\usepackage[pdftex]{hyperref} % Гиперссылки
\hypersetup{hidelinks}

% Листинги 
\usepackage{listings}

\definecolor{darkgray}{gray}{0.15}

\definecolor{teal}{rgb}{0.25,0.88,0.73}
\definecolor{gray}{rgb}{0.5,0.5,0.5}
\definecolor{b-red}{rgb}{0.88,0.25,0.41}
\definecolor{royal-blue}{rgb}{0.25,0.41,0.88}

\usepackage{listings}
\lstset{
	aboveskip=3mm,
	belowskip=3mm,
	frame=tb,
	frame=single,
	basicstyle=\footnotesize\ttfamily,
	numberstyle=\tiny\color{gray},
	keywordstyle=\color{royal-blue},
	commentstyle=\color{gray35},
	stringstyle=\color{b-red},
	numbers=left,
	numbersep=5pt,
	numberstyle=\tiny,
	showstringspaces=false, 
	captionpos=t,
	tabsize=4,
	language=Lisp
}

% какой то сложный кусок со стак эксчейндж для квадратных скобок
\makeatletter
\newenvironment{sqcases}{%
	\matrix@check\sqcases\env@sqcases
}{%
	\endarray\right.%
}
\def\env@sqcases{%
	\let\@ifnextchar\new@ifnextchar
	\left\lbrack
	\def\arraystretch{1.2}%
	\array{@{}l@{\quad}l@{}}%
}
\makeatother

% и для матриц
\makeatletter
\renewcommand*\env@matrix[1][\arraystretch]{%
	\edef\arraystretch{#1}%
	\hskip -\arraycolsep
	\let\@ifnextchar\new@ifnextchar
	\array{*\c@MaxMatrixCols c}}
\makeatother


\begin{document}

\begin{titlepage}
	
	\newgeometry{left=2cm, right=2cm, top=2.5cm, bottom=2.5cm}
	\fontsize{12pt}{12pt}\selectfont
	
	\noindent \begin{minipage}{0.13\textwidth}
		\includegraphics[width=\linewidth]{assets/bmstu-logo.png}
	\end{minipage}
	\noindent\begin{minipage}{0.85\textwidth}\centering
		\textbf{\textsc{Министерство науки и высшего образования Российской Федерации}}\\
		\textbf{\textsc{Федеральное государственное бюджетное образовательное 	учреждение высшего образования}}\\
		\textbf{\textsc{Московский государственный технический университет имени 	Н.Э.~Баумана}}\\
		\textbf{\textsc{(национальный исследовательский университет)}}\\
		\textbf{\textsc{(МГТУ им. Н.Э. Баумана)}}\\
	\end{minipage}
	
	\noindent\rule{18cm}{1.5pt}
	
	\vspace{8mm}
	
	\noindent\textnormal{ФАКУЛЬТЕТ}\hspace{5mm} \underline{\textnormal{~~~~~~~~~~~~~~~~~~«Информатика и системы управления»~~~~~~~~~~~~~~~~~~}} \newline\newline
	\textnormal{КАФЕДРА}\hspace{5mm} \underline{\textnormal{~~«Программное обеспечение ЭВМ и информационные технологии»~~}}
	\newline\newline
	\textnormal{НАПРАВЛЕНИЕ ПОДГОТОВКИ}\hspace{5mm} \underline{\textnormal{~~~~~~«09.03.04 Программная инженерия»~~~~~~~~}}
	
	\vspace{2.5cm}
	
	\begin{center}
		\Large\textbf{\textsc{ОТЧЕТ}}\\
		\Large\textbf{\textsc{ПО ЛАБОРАТОРНОЙ РАБОТЕ №18}}\\
	\end{center}
	
	\vspace{1cm}
	
	\noindent\textnormal{Название:} \hspace{15mm} \underline{\textnormal{~~~~~~Формирование и модификация списков на Prolog~~~~~~~}}\noindent
	
	\vspace{1.3cm}
	
	\noindent\textnormal{Дисциплина:} \hspace{10mm} \underline{\textnormal{~~~~~~Функциональное и логическое программирование~~~~~~}}\noindent
	
	\vspace{1.5cm}
	
	\noindent\textnormal{Студент} \hspace{17mm}
	\underline{\textnormal{{~~~~ИУ7-64Б~~~}}}
	\hspace{20mm}
	\underline{\textnormal{\hphantom{~~~~~~~~~~~~~~~~~~~~~~~~~~~}}} \hspace{14mm}
	\underline{\textnormal{~~~С. Д. Параскун~~~}}
	
	\vspace{2mm}
	\noindent\textnormal{\hphantom{Студент}} \hspace{23mm}\noindent
	\fontsize{8pt}{8pt}
	\textnormal{Группа}\hspace{40mm}\textnormal{Подпись, дата} \hspace{30mm}\noindent\textnormal{И. О. Фамилия}
	
	\vspace{0.5cm}
	
	\fontsize{12pt}{12pt}\selectfont
	\noindent\textnormal{Преподаватель} \hspace{52mm}
	\underline{\textnormal{\hphantom{~~~~~~~~~~~~~~~~~~~~~~~~~~~}}} \hspace{14mm}
	\noindent\underline{\textnormal{~~Н. Б. Толпинская~~}}
	
	\vspace{2mm}
	\noindent\textnormal{\hphantom{Студент}} \hspace{17mm}\noindent
	\fontsize{8pt}{8pt}
	\hphantom{Группа}\hspace{43mm}\textnormal{Подпись, дата} \hspace{30mm}\noindent\textnormal{И. О. Фамилия}
	
	\vspace{0.5cm}
	
	\fontsize{12pt}{12pt}\selectfont
	\noindent\textnormal{Преподаватель} \hspace{52mm}
	\underline{\textnormal{\hphantom{~~~~~~~~~~~~~~~~~~~~~~~~~~~}}} \hspace{14mm}
	\noindent\underline{\textnormal{~~Ю. В. Строганов~~}}
	
	\vspace{2mm}
	\noindent\textnormal{\hphantom{Студент}} \hspace{17mm}\noindent
	\fontsize{8pt}{8pt}
	\hphantom{Группа}\hspace{43mm}\textnormal{Подпись, дата} \hspace{30mm}\noindent\textnormal{И. О. Фамилия}
	
	\vspace{1cm}
	
	\fontsize{12pt}{12pt}\selectfont
	
	\begin{center}
		\vfill
		Москва, ~\the\year
		~г.
	\end{center}
	\restoregeometry

\end{titlepage}


\thispagestyle{empty}

\chapter{Практическое задание}

\section{Задание}

Создать базу знаний <<Собственники>>, дополнив (и минимально изменив) базу знаний, хранящую:
\begin{itemize}
	\item <<Телефонный справочник>>: Фамилия, №тел, Адрес -- структура (Город, Улица, №дома, №кв.);
	\item <<Автомобили>>: Фамилия\_владельца, Марка, Цвет, Стоимость, и др.;
	\item <<Вкладчики банков>>: Фамилия, Банк, счет, сумма, др.;
\end{itemize}
знаниями о дополнительной собственности владельца. Преобразовать знания об автомобиле к форме знаний о собственности. Виды собственности (кроме автомобиля):
\begin{itemize}
	\item Строение, стоимость и другие его характеристики;
	\item Участок, стоимость и другие его характеристики;
	\item Водный транспорт, стоимость и другие его характеристики.
\end{itemize}
Описать и использовать вариантный домен: Собственность. Владелец может иметь, но только один объект каждого вида собственности (касается и автомобиля), или не иметь некоторых видов собственности.

Используя конъюнктивное правило и разные формы задания одного вопроса (пояснить для какого № задания -- какой вопрос), обеспечить возможность поиска:
\begin{enumerate}
	\item Названий всех объектов собственности заданного субъекта;
	\item Названий и стоимости всех объектов собственности заданного субъекта;
	\item * Разработать правило, позволяющее найти суммарную стоимость всех объектов собственности заданного субъекта.
\end{enumerate}
Для 2-го пункта и одной фамилии составить таблицу, отражающую конкретный порядок работы системы, с объяснением порядка работы и особенностей использования доменов (указать конкретные T1 и T2 и полную подстановку на каждом шаге).

\section{Код программы}

\begin{lstlisting}[language=Prolog]
domains    

	surname, phone, city, street = symbol.    
	homeNumber, appartmentNumber = unsigned.
	address = address(city, street, homeNumber, appartmentNumber). 
	
	model, color = symbol.   
	price = unsigned.      
	
	bank, account = symbol.   
	sum = unsigned.    
	
	type = symbol.   
	size = unsigned.      

	ownership = building(address, price); 
	area(size, price); 
	waterTrasnport(type, color, price); 
	car(model, color, price).  

predicates    

	hasPhone(surname, phone, address).   
	hasDeposit(surname, bank, account, sum).      
	
	own(surname, ownership).   
	ownObject(surname, symbol, price).
	ownObjectPass(surname, symbol, price).  
	objectsPrice(surname, price).

clauses    
	hasPhone("Balashov", "+79741632985", address("Moscow", "Baumanskaya", 
		15, 21)).   
	hasPhone("Serov", "+79146941728", address("Lipetsk", "Gagarina", 192, 
		13)).   
	hasPhone("Paraskun", "+79172641928", address("Moscow", "Izmaylovskaya",
		73, 2)).    
	hasDeposit("Balashov", "Home-credit", "5148465849516259", 24318947).   
	hasDeposit("Balashov", "VTB", "5670148746192648", 478976).   
	hasDeposit("Paraskun", "Sberbank", "7193019871942510", 100000).
	
	own("Balashov", area(20, 139200)).   
	own("Balashov", car("BMW-Y15", "Red", 2345700)).   
	own("Paraskun", building(address("Moscow", "Lubyanka", 13, 182),
		1410000)).   
	own("Paraskun", waterTrasnport("Bike", "White", 80000)).
	
	ownObject(Surname, building, Price) :- own(Surname, building(_, Price)).
	ownObject(Surname, area, Price) :- own(Surname, area(_, Price)).
	ownObject(Surname, waterTrasnport, Price) :- own(Surname, 
		waterTrasnport(_, _, Price)).
	ownObject(Surname, car, Price) :- own(Surname, car(_, _, Price)).
	
	ownObjectPass(Surname, building, Price) :- own(Surname, building(_, 
		Price)), !.
	ownObjectPass(Surname, area, Price) :- own(Surname, area(_, Price)), !.
	ownObjectPass(Surname, waterTrasnport, Price) :- own(Surname, 
		waterTrasnport(_, _, Price)), !.
	ownObjectPass(Surname, car, Price) :- own(Surname, car(_, _, Price)), !.
	ownObjectPass(_, _, 0).
	
	objectsPrice(Surname, Price) :- 
		ownObjectPass(Surname, building, BPrice),
		ownObjectPass(Surname, area, APrice),
		ownObjectPass(Surname, waterTransport, WtPrice),
		ownObjectPass(Surname, car, CPrice),
		Price = BPrice + APrice + WtPrice + CPrice.

goal
	%ownObject("Balashov", Objects, _).
	%ownObject("Paraskun", Objects, Price).
	objectsPrice("Paraskun", Price).
\end{lstlisting}

\section{Таблицы выполнения программы}

Запрос для заданий 2:

ownObject(``Paraskun'', Objects, Price).

\begin{table}[H]
	\begin{tabular}{|p{1.2cm\small}|p{9cm\small}|p{5cm\small}|}	
		\hline
		№ шага & Сравниваемые термы; результат; подстановка, если есть & Дальнейшие действия: прямой ход или откат (к чему приводит?)\\
		\hline
		1 & Сравнение:\linebreak ownObject(``Paraskun'', Objects, Price) = hasPhone(``Balashov'', ``+79741632985'', address(``Moscow'', ``Baumanskaya'', 15, 21)). Унификация неуспешна (несовпадение функторов) & Прямой ход, переход к следующему предложению\\
		\hline
		2-10 & ... & ...\\
		\hline
		11 & Сравнение:\linebreak ownObject(``Paraskun'', Objects, Price) = ownObject(Surname, building, Price). Унификация успешна.\linebreak \textbf{Подстановка:} \{Surname=``Paraskun'', Objects=building, Price=Price\} & Новое состояние резольвенты: own(``Paraskun'', building(\_, Price))\\
		\hline
		12 & Сравнение:\linebreak own(``Paraskun'', building(\_, Price)) = hasPhone(``Balashov'', ``+79741632985'', address(``Moscow'', ``Baumanskaya'', 15, 21)). Унификация неуспешна (несовпадение функторов) & Прямой ход, переход к следующему предложению\\
		\hline
		13-19 & Унификация неуспешна & Прямой ход, переход к следующему предложению\\
		\hline
		20 & Сравнение:\linebreak own(``Paraskun'', building(\_, Price)) = own(``Paraskun'', building(address(``Moscow'', ``Lubyanka'', 13, 182), 1410000)). Унификация успешна.\linebreak \textbf{Подстановка:} \{Surname=``Paraskun'', Objects=building, Price=1410000\} & Новое состояние резольвенты: пуста\linebreak \textbf{Вывод:} Objects=building, Price=1410000\linebreak Откат, следующее предложение, \textbf{новая подстановка:} \{Surname=``Paraskun'', Objects=building, Price=Price\}\\
		\hline
		21 & Сравнение:\linebreak own(``Paraskun'', building(\_, Price)) = own(``Paraskun'', waterTrasnport(``Bike'', ``White'', 80000)). Унификация неуспешна (несовпадение термов) & Прямой ход, переход к следующему предложению\\
		\hline
		22-31 & ... (несовпадение функторов) & Откат, достижение конца БЗ, переход к следующему предложению относительно 11\\
		\hline
		32 & Сравнение:\linebreak ownObject(``Paraskun'', Objects, Price) = ownObject(Surname, area, Price). Унификация успешна.\linebreak \textbf{Подстановка:} \{Surname=``Paraskun'', Objects=area, Price=Price\} & Новое состояние резольвенты: own(``Paraskun'', area(\_, Price))\\
		\hline
		33-53 & Унификация неуспешна & Откат, достижение конца БЗ, переход к следующему предложению относительно 32\\
		\hline
	\end{tabular}
\end{table}

\begin{table}[H]
	\begin{tabular}{|p{1.2cm\small}|p{9cm\small}|p{5cm\small}|}	
		\hline
		№ шага & Сравниваемые термы; результат; подстановка, если есть & Дальнейшие действия: прямой ход или откат (к чему приводит?)\\
		\hline
		54 & Сравнение:\linebreak ownObject(``Paraskun'', Objects, Price) = ownObject(Surname, waterTransport, Price). Унификация успешна.\linebreak \textbf{Подстановка:} \{Surname=``Paraskun'', Objects=waterTransport, Price=Price\} & Новое состояние резольвенты: own(``Paraskun'', waterTransport(\_, \_, Price))\\
		\hline
		55-63 & Унификация неуспешна & Прямой ход, переход к следующему предложению\\
		\hline
		64 & Сравнение:\linebreak ownObject(``Paraskun'', Objects, Price) = own(``Paraskun'', waterTrasnport(``Bike'', ``White'', 80000)). Унификация успешна.\linebreak \textbf{Подстановка:} \{Surname=``Paraskun'', Objects=waterTransport, Price=80000\} & Новое состояние резольвенты: пуста\linebreak \textbf{Вывод:} Objects=waterTransport, Price=80000\linebreak Откат, переход к следующему предложению, \textbf{новая подстановка:} \{Surname=``Paraskun'', Objects=waterTransport, Price=Price\}\\
		\hline
		65-74 & Унификация неуспешна & Откат, достижение конца БЗ, переход к следующему предложению относительно 54\\
		\hline
		75 & Сравнение:\linebreak ownObject(``Paraskun'', Objects, Price) = ownObject(Surname, car, Price). Унификация успешна.\linebreak \textbf{Подстановка:} \{Surname=``Paraskun'', Objects=car, Price=Price\} & Новое состояние резольвенты: own(``Paraskun'', car(\_, \_, Price))\\
		\hline
		76-95 & Унификация неуспешна & Откат, достижение конца БЗ, переход к следующему предложению относительно 75\\
		\hline
		96 & Сравнение:\linebreak ownObject(``Paraskun'', Objects, Price) = ownObjectPass(Surname, building, Price).Унификация неуспешн (несовпадение функторов) & Прямой ход, переход к следующему предложению\\
		\hline
		97-101 & Унификация неуспешна & Откат, достижение конца БЗ, резольвента пуста, завершение работы\\
		\hline
	\end{tabular}
\end{table}

\end{document}